% Part: lambda-calculus
% Chapter: church-rosser
% Section: beta-reduction

\documentclass[../../../include/open-logic-section]{subfiles}

\begin{document}

\olfileid{lam}{cr}{b}

\olsection{$\beta$-తగ్గింపు}

\begin{lem}\ollabel{lem:one-par}
  $M \bredone M'$ అయితే $M \bredpar M'$.
\end{lem}
\begin{proof} $M \bredone M'$ అయితే, మూల నిరూపణ
    $M$ను $(\lambd[x][N])Q$ రూపంలో, $M'$ను
    $\Subst{N}{Q}{x}$ రూపంలో తీసుకుంటుంది; ఇక్కడ
    $x$, $N$, $Q$ తగిన చరం, పదాలు. \olref[pb]{thm:refl}
    ప్రకారం $N \bredpar N$, $Q \bredpar Q$. కనుక
    \olref[pb]{defn:bredpar}\olref[pb]{defn:bredpar4}
    ప్రకారం వెంటనే $(\lambd[x][N])Q
    \bredpar \Subst{N}{Q}{x}$.
    \sourcecorrection{OLTELAMCRB-001}{మూల వాక్యం ప్రతి $M \bredone M'$ను పై మూలస్థాన రెడెక్స్ రూపంలోనే తీసుకుంటుంది. కానీ ముందరి $\beta$-సంకోచన నిర్వచనం $\bredone$ను అనుకూల సంబంధంగా ఇస్తుంది; అమూర్తీకరణ లేదా ప్రయోగం లోపల కూడా ఒక దశ సంకోచించవచ్చు. ఇక్కడి గణన మూలస్థాన సందర్భాన్నే నిరూపిస్తుంది; మిగతా సందర్భాలకు ఆగమన వాదనను మూలం ఇవ్వలేదు. ఉపసిద్ధాంతం నిరూపణ పూర్తయిందని ఈ అనువాదం ప్రకటించదు.}
\end{proof}

\begin{lem}\ollabel{lem:par-red}
  $M \bredpar M'$ అయితే $M \bred M'$.
\end{lem}

\begin{proof} $M \bredpar M'$ యొక్క
  \tetoken{వ్యుత్పత్తి}{!!{derivation}}పై ఆగమన పద్ధతిలో నిరూపిస్తాం.
  \begin{enumerate}
  \item చివరి నియమం \olref[pb]{defn:bredpar1}: అప్పుడు
    $M$, $M'$ రెండూ $x$ మాత్రమే; $x \bred x$.
  \item చివరి నియమం \olref[pb]{defn:bredpar2}: $M$ అనేది
    $\lambd[x][N]$, $M'$ అనేది $\lambd[x][N']$;
    ఇక్కడ $x$, $N$, $N'$ తగిన చరం, పదాలు;
    $N \bredpar N'$. ఆగమన పరికల్పన ప్రకారం
    $N \bred N'$. కనుక $N \bred N'$కు వాడిన
    $\bredone$-సంకోచన క్రమాన్నే అమూర్తీకరణ లోపల
    వర్తింపజేసి $\lambd[x][N] \bred
    \lambd[x][N']$ పొందుతాం.
  \item చివరి నియమం \olref[pb]{defn:bredpar3}: $M$ అనేది
    $PQ$, $M'$ అనేది $P'Q'$; ఇక్కడ $P$, $Q$,
    $P'$, $Q'$ తగిన పదాలు; $P \bredpar P'$,
    $Q \bredpar Q'$. ఆగమన పరికల్పన ప్రకారం
    $P \bred P'$, $Q \bred Q'$. ముందు $P \bred P'$
    క్రమాన్నీ, తరువాత $Q \bred Q'$ క్రమాన్నీ ప్రయోగంలో
    వర్తింపజేస్తే $PQ \bred P'Q'$.
  \item చివరి నియమం \olref[pb]{defn:bredpar4}: $M$ అనేది
    $(\lambd[x][N])Q$, $M'$ అనేది $\Subst{N'}{Q'}{x}$;
    ఇక్కడ $x$, $N$, $N'$, $Q$, $Q'$ తగిన చరం,
    పదాలు; $N \bredpar N'$, $Q \bredpar Q'$.
    \sourcecorrection{OLTELAMCRB-002}{మూలంలోని ఈ పదాల జాబితా $x,N,M',Q,Q'$ అని ఉంది; అదే వాక్యంలో $M'=\Subst{N'}{Q'}{x}$, తరువాతి పరికల్పనలో $N \bredpar N'$ అని ఉన్నాయి. అవసరమైన శరీరపు ఫలిత పదం $N'$ను జాబితాలో పునరుద్ధరించాం; ఇతర సూత్రాలు మార్చలేదు.}
    ఆగమన పరికల్పన ప్రకారం $Q \bred Q'$,
    $N \bred N'$. కాబట్టి ముందుగా $N \bred N'$,
    తరువాత $Q \bred Q'$ తగ్గింపులను చేసి, చివరికి
    $(\lambd[x][N'])Q'$ను $\Subst{N'}{Q'}{x}$కు
    సంకోచిస్తే $(\lambd[x][N])Q \bred
    \Subst{N'}{Q'}{x}$ వస్తుంది.
  \end{enumerate}
\end{proof}

\begin{lem}\ollabel{lem:str}
  $\bredpar$ను కలిగి ఉండే కనిష్ఠ సంక్రమణ సంబంధం $\bred$.
\end{lem}

\begin{proof}
  $\bredpar$ను కలిగి ఉండే కనిష్ఠ సంక్రమణ సంబంధాన్ని
  $\xred$ అనుకుందాం.

  $\bred \subseteq \xred$: $M \bred M'$ అని అనుకుందాం;
  అంటే $M \ident M_1
  \bredone \dots \bredone M_k \ident M'$.
  \olref{lem:one-par} ప్రకారం $M
  \ident M_1 \bredpar \dots \bredpar M_k \ident M'$.
  $\xred$లో $\bredpar$ ఉంది; అది సంక్రమణ సంబంధం
  కూడా. కాబట్టి $M \xred M'$.

  $\xred \subseteq \bred$: $M \xred M'$ అని అనుకుందాం;
  అంటే $M
  \ident M_1 \bredpar \dots \bredpar M_k \ident M'$.
  \olref{lem:par-red} ప్రకారం $M \ident M_1
  \bred \dots \bred M_k \ident M'$. $\bred$
  సంక్రమణ సంబంధం కనుక $M \bred M'$.
\end{proof}

\begin{thm}\ollabel{thm:cr}
  $\bred$ చర్చ్--రోసర్ లక్షణాన్ని తీరుస్తుంది.
\end{thm}

\begin{proof}
  \olref[dap]{thm:str}, \olref[pb]{thm:cr},
  \olref{lem:str} నుంచి నేరుగా వస్తుంది.
  \sourcecorrection{OLTELAMCRB-003}{ఈ తుది వాదన ముందరి సమాంతర-తగ్గింపు చర్చ్--రోసర్ సిద్ధాంతంపైనా, ఇక్కడి \olref{lem:one-par}పైనా ఆధారపడుతుంది. ముందరి OLTELAMCRPB-003 ప్రతిస్థాపన నిరూపణ ఖాళీ, ఇక్కడి OLTELAMCRB-001 అనుకూల-సందర్భ ఖాళీ ఇంకా పూరించలేదు. మూల ఆధార సూచనను అలాగే ఉంచాం; దీన్ని స్వతంత్రంగా ధ్రువీకరించిన పూర్తి నిరూపణగా పేర్కొనడం లేదు.}
\end{proof}

\end{document}
